% =============================================================
% Resúmenes en castellano, valenciano e inglés.
% Los tres son obligatorios según la plantilla de la EPSG.
% Cada uno tiene entre 250 y 830 palabras.
% =============================================================

\chapter*{Resumen}

Este Trabajo de Fin de Grado desarrolla una metaheurística capaz de generar mazos competitivos para el formato Estándar de Magic: The Gathering. El juego presenta un metajuego dinámico, con cartas que rotan cada pocos meses y arquetipos dominantes que cambian con frecuencia. El sistema construye mazos que pueden enfrentarse a esos arquetipos manteniendo coherencia interna y respetando las reglas del formato. El espacio de búsqueda incluye alrededor de cuatro mil cartas con restricciones duras como el máximo de cuatro copias por carta, los 60 cartas exactas por mazo y la consistencia cromática del manabase.

La solución implementada es un algoritmo genético con dos características que lo distinguen de un enfoque clásico. Primera: los operadores de mutación y cruce trabajan a nivel de \emph{pack} de cuatro copias, replicando cómo construye un jugador humano. Decidir si una carta entra al mazo equivale a decidir si entran cuatro copias, no una. Segunda: el fitness combina el rendimiento real en partidas simuladas con el motor Forge (al 70\%) y una métrica de calidad estructural sensible al arquetipo detectado (al 30\%). La población evoluciona durante decenas de generaciones bajo presión selectiva tipo Swiss Tournament, con elitismo soportado por un Hall of Fame con cuotas por arquetipo.

Los resultados confirman la viabilidad del enfoque. En la comparativa frente al baseline sin operadores a nivel de pack, la nueva aproximación reduce los \emph{singletons} sueltos de 3,55 a 0,05 por mazo (una bajada del 95\%) y consolida 7,88 series de cuatro copias por mazo. El run final, de 43 generaciones sobre una población de 40 mazos, alcanza un pico de fitness de 0,9755 en la generación 39, lo que sitúa el mejor mazo al 93\% del recorrido evolutivo. Esa progresión escalonada descarta la convergencia prematura típica de los algoritmos genéticos mal calibrados. El Hall of Fame final reúne mazos con cartas reconocidas del meta competitivo actual de Standard, como Scalestorm Summoner, Leyline Axe o el paquete de quemado Boros.

El trabajo también identifica límites del enfoque actual. Cada ejecución del algoritmo converge a un único par cromático, por lo que recorrer el ecosistema arquetípico al completo exigiría varias ejecuciones con semillas distintas. Una segunda exploración, basada en evaluar la población contra mazos tier-1 del meta real, se descartó tras detectar un sesgo medible del simulador frente a arquetipos de juego complejo. La memoria documenta ambos hallazgos y propone líneas concretas de continuación.

\textbf{Palabras clave:} algoritmo genético, Magic: The Gathering, generación de mazos, metaheurística, simulación.

\chapter*{Resum}
\selectlanguage{catalan}
% TODO: traducir el resumen al valenciano.

\textbf{Paraules clau:}
% TODO: palabras clave en valenciano.

\selectlanguage{english}
\chapter*{Abstract}

This Bachelor's Thesis develops a metaheuristic capable of generating competitive decks for the Standard format of Magic: The Gathering. The game features a dynamic metagame, with cards rotating every few months and dominant archetypes shifting frequently. The system builds decks that can face those archetypes while keeping internal coherence and complying with the format rules. The search space contains around four thousand cards under hard constraints such as the maximum of four copies per card, the requirement of exactly 60 cards per deck and the chromatic consistency of the manabase.

The implemented solution is a genetic algorithm with two features that set it apart from a classical approach. First: the mutation and crossover operators work at the level of a \emph{pack} of four copies, mirroring how a human player builds. Deciding whether a card enters the deck is equivalent to deciding whether four copies enter, not one. Second: the fitness combines real performance from simulated matches in the Forge engine (at 70\%) with an archetype-aware structural quality metric (at 30\%). The population evolves over dozens of generations under Swiss Tournament selection pressure, with elitism backed by a Hall of Fame that holds quotas per archetype.

The results confirm the viability of the approach. Compared to a baseline without pack-level operators, the new design reduces stray \emph{singletons} from 3.55 to 0.05 per deck (a 95\% drop) and consolidates 7.88 four-copy sets per deck. The final run, 43 generations over a population of 40 decks, reaches a fitness peak of 0.9755 at generation 39, which places the best deck at 93\% of the evolutionary path. That stepwise progression rules out the early convergence typical of poorly calibrated genetic algorithms. The final Hall of Fame gathers decks with cards from the current competitive Standard meta, such as Scalestorm Summoner, Leyline Axe or the Boros burn package.

The work also flags limits of the current approach. Each run of the algorithm converges to a single colour pair, so covering the full archetype ecosystem would demand several runs with different seeds. A second exploration, based on evaluating the population against tier-1 decks from the actual meta, was dropped after detecting a measurable simulator bias against archetypes with complex play. The thesis documents both findings and proposes concrete continuation lines.

\textbf{Keywords:} genetic algorithm, Magic: The Gathering, deck generation, metaheuristic, simulation.

\selectlanguage{spanish}
